\section{辩护制度}
\begin{enumerate}
  \item 围绕保障嫌疑人/被告人的辩护权作为前提、核心来进行制度的展开和设计。
  \item 辩护、辩护权、辩护制度。
  \item 有效辩护原则。形式 + 实质、\zd{整个诉讼过程都重要}、请有资格有能力的人、法律援助制度。
  \item 意义。发现真相、保障程序正义、法制宣传教育。
  \item 分类。
        \begin{enumerate}
          \item 自行辩护。随时随地，想辩就辩。
          \item 委托辩护。
                \begin{enumerate}
                  \item \zd{告知义务。自第一次讯问、检察院收到审查起诉材料 3 日内、法院受案 3 日内}。
                  \item 公诉。律师 \yjt 自\zd{第一次讯问 \hz 采取强制措施之日}起。普通公民 \yjt \zd{移送审查起诉之日}起。
                  \item 在押。监护人、近亲属可以代为委托。
                  \item 自诉。随时。
                \end{enumerate}
          \item 法律援助辩护。
                \begin{enumerate}
                  \item 前提是没有辩护人 \yjt 辩护人到位则退出。只能是律师。
                  \item 方式。
                        \begin{enumerate}
                          \item 申请。穷。
                          \item 强制。\bs{盲聋哑、半疯傻、死无缺、未长大。应当通知法援机构}。
                          \item 裁量。
                        \end{enumerate}
                \end{enumerate}
          \item 值班律师。\zd{替补}、填坑的。目前阶段的功能是证明公权行为的正确性。哈哈哈哈。
                \begin{enumerate}
                  \item 没有委托辩护人 \hz 法援没派律师。
                  \item 辩护制度的重要补充（\zd{\bdy 辩护权}）、法律援助制度的重要组成部分、\bs{身份不是辩护人}。
                  \item 职责。\zd{提供咨询、出}量刑\zd{建议}；\zd{申请变更}强制措施、\zd{提意见；申请法援或其他}。
                  \item 权利。\zd{会见、阅卷、意见}。类推辩护律师。\zd{没有出庭、取证的权利}。
                  \item \bs{无期死刑有经验（3 年），未成年人懂特点}。新增。
                  \item 不可转为辩护律师。公益法律服务变成了案源拓展活动。
                \end{enumerate}
        \end{enumerate}
  \item \kd{辩护人}。
        \begin{enumerate}
          \item 同一个案件不管时间，\bs{1 人最多找 2 辩；1 人只为 1 人辩}。\al{甲乙犯罪，乙在逃，丙担任甲辩护人，判决生效 3 年后乙归案，丙不能担任乙辩护人}。
          \item 范围（能找谁）。积极，消极。
                \begin{enumerate}
                  \item 绝对禁止。\bs{刑罚执行中、自由受限、能力有限}。
                  \item 相对禁止。\zd{公检法安监现}职人员，\zd{陪}审员，\zd{利}害关系，\zd{外}国人、被\zd{开}除\zd{公}职或\zd{吊}销律师\zd{执照}。
                  \item 法检限制。\zd{2 年终身家庭店}。不能用律师身份。
                  \item \sjt 例外。除非他是近亲属或监护人。
                \end{enumerate}
          \item 地位独立和职责辩护，但\zd{不能够超出指控的罪名}。\al{他不仅犯了故意杀人罪，还主张他有强奸罪}。
          \item 辩护权。归属于嫌疑人/被告人及其辩护人。\bs{做题先区分身份，非律辩护矮半截}。
                \begin{enumerate}
                  \item 不同的阶段权限不同，\zd{侦查阶段权限最小}。
                  \item \bs{侦查阶段。无阅卷、无核证、无非律}。
                        \begin{enumerate}
                          \item 会见。\bs{非律会见要许可，律师一般不用批。律师会见凭 3 证，危恐除外要许可}。
                                \begin{itemize}
                                  \item \zd{不能监听}，可以监视。
                                  \item 3 证。律师证、事务所证明、委托书 \hz 法律援助公函。
                                  \item 可以\zd{带 1 名律师助理}。
                                  \item 会见应当遵守看管场所规定；遵守法庭秩序；依法取证；不得违规会见、贿赂司法人员。
                                \end{itemize}
                          \item 通信。\bs{国公人身串毁证，信件可被截复删}。国家、公共安全、人身危险、串供、毁灭证据。
                          \item 了解案情、提供咨询、提供帮助。
                          \item 代理申诉、控告。同级（公法侵害）\hz 上级检察院。
                          \item 申请变更、解除强制措施。\bs{当}事人、\bs{法}定代理人、\bs{近}亲属、\bs{辩}护人。
                          \item 主动听。\bs{批捕未成年，审查起诉中，二审不开庭，死刑上抗案，认罪又认罚}。
                        \end{enumerate}
                  \item 阅卷。\bs{非律阅卷要许可，律师阅卷不用批，均自审查起诉起}。
                        \begin{itemize}
                          \item 扩大解释。阅 \yjt 复印拷贝。卷 \yjt 案件材料、证据等。
                          \item 审查起诉后、地点为检察院或者法院、\zd{律师不收费}。
                        \end{itemize}
                  \item 取证。\bs{非律没有取证权，律师取证需同意。控方证人双重许，若有困难找法检}。
                  \item 知情权。重要的诉讼节点都要告知 \zd{And(嫌疑人，辩护律师 \hz 辩护人)}。
                  \item 参加法庭调查、辩论。值班律师他只是个替补，不可以。
                  \item 人身保障。
                  \item 拒绝辩护。
                \end{enumerate}
          \item 义务。
                \begin{enumerate}
                  \item 不得干扰司法机关诉讼活动。及时告知接受委托的情况。
                  \item 有利证据及时告知。\zd{不在场、不够大、不正常}。
                  \item 保密、揭发违法行为。\bs{过去统统要保密，将来国、公、人身要揭发}。
                \end{enumerate}
          \item 辩护律师涉嫌犯罪。
                \begin{enumerate}
                  \item 干扰诉讼活动行为 + 可能追究刑事责任 \yjt 上一级 \hz 上一级指定。
                  \item 原来的管可能\zd{涉嫌打击报复}。
                \end{enumerate}
          \item 辩护意见。无罪、有罪；量刑可当庭 \hz 庭后。
        \end{enumerate}
  \item 被告人拒绝辩护。
        \begin{enumerate}
          \item 强制辩护。要\zd{有理}由 + \zd{1 次}机会 + 必须\zd{有人辩}护。
          \item 其他辩护。不用\zd{（无）理}由 + \zd{2 次}机会 + 最终只能\zd{自己辩}。
        \end{enumerate}
  \item 辩护的内容。是否构成犯\zd{罪}、承担刑事\zd{责}任、量\zd{刑}、罪\zd{名}、\zd{过}错、\zd{主从}犯/\zd{分析证据和程序}。
  \item 辩护的分类。无罪、罪名、罪数、量刑、程序。主观题。
\end{enumerate}
